% Compile with: uplatex file.tex && dvipdfmx file.dvi
%           or: latexmk -uplatex file.tex
%\documentclass[uplatex,dvipdfmx]{jsarticle}
\documentclass[dvipdfmx,a4paper,11pt]{jsarticle}
\usepackage{amsmath, amssymb}
\usepackage{tikz}

\title{scenario\_ac\_impedance}
\date{\today}

\begin{document}
\maketitle

直列RLC回路のインピーダンスを扱うため、抵抗 R・インダクタンス L・角周波数 w（$\omega$）・キャパシタンス C を実数のスカラーパラメータとして準備する。

R、L、w、C をスカラーとする

\[
  R, L, w, C
\]

抵抗R・インダクタL・キャパシタCを直列接続した回路のインピーダンスは、複素数 Z=R+j$\omega$L+1/(j$\omega$C) で与えられる（j は虚数単位。この electrical\_engineering ドメインでは工学の慣習に従い j を虚数単位、i を電流として扱う）。

$Z = R + j w L + 1/(j w C)$

\begin{equation}
  Z = R + j w L + 1/(j w C)
  \label{eq:E1}
\end{equation}

インピーダンス Z の実部（抵抗成分）を取り出す。

E1 の実部を求めて

\begin{equation}
  R
  \label{eq:E2}
\end{equation}

インピーダンス Z の虚部（リアクタンス成分）を取り出す。エネルギーを蓄えるL・Cによって生じる、電流と電圧の位相のずれの大きさを表す。

E1 の虚部を求めて

\begin{equation}
  L w-\frac{1}{C w}
  \label{eq:E3}
\end{equation}

共振とは、リアクタンス（虚部）がゼロになる状態を指す。ここでは虚部の値を0と置いた等式を作る。

$E3 = 0$

\begin{equation}
  L w-1/(C w) = 0
  \label{eq:E4}
\end{equation}

この等式を角周波数 w について解くと、共振の起きる角周波数が得られる。符号違いの2つの解のうち、物理的に意味を持つのは正の解 $\omega$0=1/$\surd$(LC) である。

E4 を w について解いて

\begin{equation}
  w = -\sqrt{\frac{1}{C L}}
  \label{eq:E5}
\end{equation}

\begin{equation}
  w = \sqrt{\frac{1}{C L}}
  \label{eq:E6}
\end{equation}

\end{document}